% Options for packages loaded elsewhere
% Options for packages loaded elsewhere
\PassOptionsToPackage{unicode}{hyperref}
\PassOptionsToPackage{hyphens}{url}
\PassOptionsToPackage{dvipsnames,svgnames,x11names}{xcolor}
%
\documentclass[
  letterpaper,
  DIV=11,
  numbers=noendperiod]{scrartcl}
\usepackage{xcolor}
\usepackage{amsmath,amssymb}
\setcounter{secnumdepth}{-\maxdimen} % remove section numbering
\usepackage{iftex}
\ifPDFTeX
  \usepackage[T1]{fontenc}
  \usepackage[utf8]{inputenc}
  \usepackage{textcomp} % provide euro and other symbols
\else % if luatex or xetex
  \usepackage{unicode-math} % this also loads fontspec
  \defaultfontfeatures{Scale=MatchLowercase}
  \defaultfontfeatures[\rmfamily]{Ligatures=TeX,Scale=1}
\fi
\usepackage{lmodern}
\ifPDFTeX\else
  % xetex/luatex font selection
\fi
% Use upquote if available, for straight quotes in verbatim environments
\IfFileExists{upquote.sty}{\usepackage{upquote}}{}
\IfFileExists{microtype.sty}{% use microtype if available
  \usepackage[]{microtype}
  \UseMicrotypeSet[protrusion]{basicmath} % disable protrusion for tt fonts
}{}
\makeatletter
\@ifundefined{KOMAClassName}{% if non-KOMA class
  \IfFileExists{parskip.sty}{%
    \usepackage{parskip}
  }{% else
    \setlength{\parindent}{0pt}
    \setlength{\parskip}{6pt plus 2pt minus 1pt}}
}{% if KOMA class
  \KOMAoptions{parskip=half}}
\makeatother
% Make \paragraph and \subparagraph free-standing
\makeatletter
\ifx\paragraph\undefined\else
  \let\oldparagraph\paragraph
  \renewcommand{\paragraph}{
    \@ifstar
      \xxxParagraphStar
      \xxxParagraphNoStar
  }
  \newcommand{\xxxParagraphStar}[1]{\oldparagraph*{#1}\mbox{}}
  \newcommand{\xxxParagraphNoStar}[1]{\oldparagraph{#1}\mbox{}}
\fi
\ifx\subparagraph\undefined\else
  \let\oldsubparagraph\subparagraph
  \renewcommand{\subparagraph}{
    \@ifstar
      \xxxSubParagraphStar
      \xxxSubParagraphNoStar
  }
  \newcommand{\xxxSubParagraphStar}[1]{\oldsubparagraph*{#1}\mbox{}}
  \newcommand{\xxxSubParagraphNoStar}[1]{\oldsubparagraph{#1}\mbox{}}
\fi
\makeatother

\usepackage{color}
\usepackage{fancyvrb}
\newcommand{\VerbBar}{|}
\newcommand{\VERB}{\Verb[commandchars=\\\{\}]}
\DefineVerbatimEnvironment{Highlighting}{Verbatim}{commandchars=\\\{\}}
% Add ',fontsize=\small' for more characters per line
\usepackage{framed}
\definecolor{shadecolor}{RGB}{241,243,245}
\newenvironment{Shaded}{\begin{snugshade}}{\end{snugshade}}
\newcommand{\AlertTok}[1]{\textcolor[rgb]{0.68,0.00,0.00}{#1}}
\newcommand{\AnnotationTok}[1]{\textcolor[rgb]{0.37,0.37,0.37}{#1}}
\newcommand{\AttributeTok}[1]{\textcolor[rgb]{0.40,0.45,0.13}{#1}}
\newcommand{\BaseNTok}[1]{\textcolor[rgb]{0.68,0.00,0.00}{#1}}
\newcommand{\BuiltInTok}[1]{\textcolor[rgb]{0.00,0.23,0.31}{#1}}
\newcommand{\CharTok}[1]{\textcolor[rgb]{0.13,0.47,0.30}{#1}}
\newcommand{\CommentTok}[1]{\textcolor[rgb]{0.37,0.37,0.37}{#1}}
\newcommand{\CommentVarTok}[1]{\textcolor[rgb]{0.37,0.37,0.37}{\textit{#1}}}
\newcommand{\ConstantTok}[1]{\textcolor[rgb]{0.56,0.35,0.01}{#1}}
\newcommand{\ControlFlowTok}[1]{\textcolor[rgb]{0.00,0.23,0.31}{\textbf{#1}}}
\newcommand{\DataTypeTok}[1]{\textcolor[rgb]{0.68,0.00,0.00}{#1}}
\newcommand{\DecValTok}[1]{\textcolor[rgb]{0.68,0.00,0.00}{#1}}
\newcommand{\DocumentationTok}[1]{\textcolor[rgb]{0.37,0.37,0.37}{\textit{#1}}}
\newcommand{\ErrorTok}[1]{\textcolor[rgb]{0.68,0.00,0.00}{#1}}
\newcommand{\ExtensionTok}[1]{\textcolor[rgb]{0.00,0.23,0.31}{#1}}
\newcommand{\FloatTok}[1]{\textcolor[rgb]{0.68,0.00,0.00}{#1}}
\newcommand{\FunctionTok}[1]{\textcolor[rgb]{0.28,0.35,0.67}{#1}}
\newcommand{\ImportTok}[1]{\textcolor[rgb]{0.00,0.46,0.62}{#1}}
\newcommand{\InformationTok}[1]{\textcolor[rgb]{0.37,0.37,0.37}{#1}}
\newcommand{\KeywordTok}[1]{\textcolor[rgb]{0.00,0.23,0.31}{\textbf{#1}}}
\newcommand{\NormalTok}[1]{\textcolor[rgb]{0.00,0.23,0.31}{#1}}
\newcommand{\OperatorTok}[1]{\textcolor[rgb]{0.37,0.37,0.37}{#1}}
\newcommand{\OtherTok}[1]{\textcolor[rgb]{0.00,0.23,0.31}{#1}}
\newcommand{\PreprocessorTok}[1]{\textcolor[rgb]{0.68,0.00,0.00}{#1}}
\newcommand{\RegionMarkerTok}[1]{\textcolor[rgb]{0.00,0.23,0.31}{#1}}
\newcommand{\SpecialCharTok}[1]{\textcolor[rgb]{0.37,0.37,0.37}{#1}}
\newcommand{\SpecialStringTok}[1]{\textcolor[rgb]{0.13,0.47,0.30}{#1}}
\newcommand{\StringTok}[1]{\textcolor[rgb]{0.13,0.47,0.30}{#1}}
\newcommand{\VariableTok}[1]{\textcolor[rgb]{0.07,0.07,0.07}{#1}}
\newcommand{\VerbatimStringTok}[1]{\textcolor[rgb]{0.13,0.47,0.30}{#1}}
\newcommand{\WarningTok}[1]{\textcolor[rgb]{0.37,0.37,0.37}{\textit{#1}}}

\usepackage{longtable,booktabs,array}
\newcounter{none} % for unnumbered tables
\usepackage{calc} % for calculating minipage widths
% Correct order of tables after \paragraph or \subparagraph
\usepackage{etoolbox}
\makeatletter
\patchcmd\longtable{\par}{\if@noskipsec\mbox{}\fi\par}{}{}
\makeatother
% Allow footnotes in longtable head/foot
\IfFileExists{footnotehyper.sty}{\usepackage{footnotehyper}}{\usepackage{footnote}}
\makesavenoteenv{longtable}
\usepackage{graphicx}
\makeatletter
\newsavebox\pandoc@box
\newcommand*\pandocbounded[1]{% scales image to fit in text height/width
  \sbox\pandoc@box{#1}%
  \Gscale@div\@tempa{\textheight}{\dimexpr\ht\pandoc@box+\dp\pandoc@box\relax}%
  \Gscale@div\@tempb{\linewidth}{\wd\pandoc@box}%
  \ifdim\@tempb\p@<\@tempa\p@\let\@tempa\@tempb\fi% select the smaller of both
  \ifdim\@tempa\p@<\p@\scalebox{\@tempa}{\usebox\pandoc@box}%
  \else\usebox{\pandoc@box}%
  \fi%
}
% Set default figure placement to htbp
\def\fps@figure{htbp}
\makeatother





\setlength{\emergencystretch}{3em} % prevent overfull lines

\providecommand{\tightlist}{%
  \setlength{\itemsep}{0pt}\setlength{\parskip}{0pt}}



 


\KOMAoption{captions}{tableheading}
\makeatletter
\@ifpackageloaded{tcolorbox}{}{\usepackage[skins,breakable]{tcolorbox}}
\@ifpackageloaded{fontawesome5}{}{\usepackage{fontawesome5}}
\definecolor{quarto-callout-color}{HTML}{909090}
\definecolor{quarto-callout-note-color}{HTML}{0758E5}
\definecolor{quarto-callout-important-color}{HTML}{CC1914}
\definecolor{quarto-callout-warning-color}{HTML}{EB9113}
\definecolor{quarto-callout-tip-color}{HTML}{00A047}
\definecolor{quarto-callout-caution-color}{HTML}{FC5300}
\definecolor{quarto-callout-color-frame}{HTML}{acacac}
\definecolor{quarto-callout-note-color-frame}{HTML}{4582ec}
\definecolor{quarto-callout-important-color-frame}{HTML}{d9534f}
\definecolor{quarto-callout-warning-color-frame}{HTML}{f0ad4e}
\definecolor{quarto-callout-tip-color-frame}{HTML}{02b875}
\definecolor{quarto-callout-caution-color-frame}{HTML}{fd7e14}
\makeatother
\makeatletter
\@ifpackageloaded{caption}{}{\usepackage{caption}}
\AtBeginDocument{%
\ifdefined\contentsname
  \renewcommand*\contentsname{Table of contents}
\else
  \newcommand\contentsname{Table of contents}
\fi
\ifdefined\listfigurename
  \renewcommand*\listfigurename{List of Figures}
\else
  \newcommand\listfigurename{List of Figures}
\fi
\ifdefined\listtablename
  \renewcommand*\listtablename{List of Tables}
\else
  \newcommand\listtablename{List of Tables}
\fi
\ifdefined\figurename
  \renewcommand*\figurename{Figure}
\else
  \newcommand\figurename{Figure}
\fi
\ifdefined\tablename
  \renewcommand*\tablename{Table}
\else
  \newcommand\tablename{Table}
\fi
}
\@ifpackageloaded{float}{}{\usepackage{float}}
\floatstyle{ruled}
\@ifundefined{c@chapter}{\newfloat{codelisting}{h}{lop}}{\newfloat{codelisting}{h}{lop}[chapter]}
\floatname{codelisting}{Listing}
\newcommand*\listoflistings{\listof{codelisting}{List of Listings}}
\makeatother
\makeatletter
\makeatother
\makeatletter
\@ifpackageloaded{caption}{}{\usepackage{caption}}
\@ifpackageloaded{subcaption}{}{\usepackage{subcaption}}
\makeatother
\makeatletter
\@ifpackageloaded{fontawesome5}{}{\usepackage{fontawesome5}}
\makeatother
\usepackage{bookmark}
\IfFileExists{xurl.sty}{\usepackage{xurl}}{} % add URL line breaks if available
\urlstyle{same}
\hypersetup{
  pdftitle={Introduction to High Performance Computing},
  colorlinks=true,
  linkcolor={blue},
  filecolor={Maroon},
  citecolor={Blue},
  urlcolor={Blue},
  pdfcreator={LaTeX via pandoc}}


\title{Introduction to High Performance Computing}
\usepackage{etoolbox}
\makeatletter
\providecommand{\subtitle}[1]{% add subtitle to \maketitle
  \apptocmd{\@title}{\par {\large #1 \par}}{}{}
}
\makeatother
\subtitle{ICCS Summer School 2026}
\author{Chris Edsall \and Tom Meltzer}
\date{}
\begin{document}
\maketitle


\subsection{High Performance
Computing}\label{high-performance-computing}

Working definition:

\begin{quote}
A computing resource that is larger than can be provided by one laptop
or server
\end{quote}

\subsection{Supercomputers and
clusters}\label{supercomputers-and-clusters}

\subsubsection{Supercomputer}\label{supercomputer}

One of the most performant computers in the world at a particular point
in time.

\subsubsection{Cluster}\label{cluster}

An architecture for combining a number of servers, storage and
networking to act on concert.

Most supercomputers for the past few decades have been clusters.

\subsection{Applications of HPC}\label{applications-of-hpc}

Why would I need a supercomputer?

Three traditional applications:

\begin{itemize}
\tightlist
\item
  Nuclear
\item
  Chemical
\item
  Climate / weather
\end{itemize}

Now, AI

\subsection{Floating Point}\label{floating-point}

Computer maths is not people maths

\begin{verbatim}
>>> 0.1 + 0.2
\end{verbatim}

\subsection{Floating Point}\label{floating-point-1}

\begin{verbatim}
>>> 0.1 + 0.2
0.30000000000000004
\end{verbatim}

\subsection{Floating Point}\label{floating-point-2}

\begin{itemize}
\tightlist
\item
  '60s and '70s many vendor implementations
\item
  Standardised in 1982 as IEEE 754
\end{itemize}

\subsection{FLOPS}\label{flops}

One FLOPS == one floating point operation per second.

\begin{itemize}
\tightlist
\item
  TF terraflops
\item
  PF petaflops
\item
  EF exaflops
\end{itemize}

Conventionally these are 64-bit (``double precision'') FLOPS

\subsection{``AI'' FLOPS}\label{ai-flops}

\begin{itemize}
\tightlist
\item
  smaller data formats

  \begin{itemize}
  \tightlist
  \item
    float16
  \item
    bfloat16
  \item
    int8
  \end{itemize}
\end{itemize}

\subsection{FLOPS}\label{flops-1}

\includegraphics[width=0.65\linewidth,height=\textheight,keepaspectratio]{imgs/what-kind-flops-goose.jpg}

Image source:
\href{https://bsky.app/profile/fclc.bsky.social/post/3lc4qpte3ys2o}{Felix
LeClair}

\subsection{Benchmarks}\label{benchmarks}

A benchmark is a particular known and specified workload which can be
repeated on different systems and the performance compared.

A typical weather related one is WRF running the CONUS 2.5km
configuration.

\subsection{HPL}\label{hpl}

LINPACK is a software library for performing numerical linear algebra

LINPACK makes use of the BLAS (Basic Linear Algebra Subprograms)
libraries for performing basic vector and matrix operations.

The LINPACK benchmarks appeared initially as part of the LINPACK user's
manual. The parallel LINPACK benchmark implementation called HPL (High
Performance Linpack) is used to benchmark and rank supercomputers for
the TOP500 list.

\subsection{Top500 list}\label{top500-list}

\subsection{Exercise 1}\label{exercise-1}

Got to the Top500 site at https://top500.org/

\begin{itemize}
\tightlist
\item
  Use the sublist generator to find the largest HPC systems in your
  country
\item
  What is the ratio of Rmax performance between the number 1 system in
  June 1993 and the June 2026 number 1
\end{itemize}

\subsection{Cluster architecture}\label{cluster-architecture}

Before we get to the computing infrastructure there is the underpinning
building and plant (power, cooling) required

\subsection{Cluster architecture}\label{cluster-architecture-1}

\pandocbounded{\includegraphics[keepaspectratio]{imgs/Cluster.drawio.png}}

\subsection{Nodes}\label{nodes}

The name comes from the terminology of mathematical graphs - nodes and
edges.

You can think of a node as a single server - one computer that an
instance of an operating system

\subsection{Login Nodes}\label{login-nodes}

These are your entry point on to the cluster

Usually accessable from the outside world.

Often more than one (sometimes multiple login nodes use the same DNS
name, e.g . \texttt{login.hpc.cam.ac.uk})

Shared with multiple users.

DO NOT RUN COMPUTE JOBS ON THE LOGIN NODE

\subsection{Compute nodes}\label{compute-nodes}

These are the nodes that do the heavy lifting computing work.

Normally managed by the job scheduler - you don't usually log in to them
directly.

Quite often for the exclusive use of one user for the duration of their
job.

N.B. On some clusters compute nodes can be of a different architecture
to the login nodes.

\subsection{Shared storage}\label{shared-storage}

Compute nodes sometimes have on node disk storage.

Ther is normally some large storage that is visible to all the compute
nodes.

Since this is a shared resource an anti-social user can affect the
performance of other users.

\subsection{Interconnect}\label{interconnect}

Connects the compute nodes, login nodes and storage

Usually faster (higher bandwidth, lower latency) than comoddity ethernet
networking.

It's what makes a supercomputer super.

examples:

\begin{itemize}
\tightlist
\item
  Infiniband
\item
  Omnipath
\item
  Slingshot
\end{itemize}

\subsection{Connecting to CSD3}\label{connecting-to-csd3}

\begin{itemize}
\tightlist
\item
  SSH

  \begin{itemize}
  \tightlist
  \item
    linux / mac should be built in
  \item
    Windows has openssh nowadays
  \item
    PuTTY
  \item
    MobaXTerm
  \end{itemize}
\item
  \href{https://carpentries-incubator.github.io/hpc-intro/11-connecting/index.html}{HPC
  Carpentry advice on connecting}
\item
  \href{https://chrome.google.com/webstore/detail/iodihamcpbpeioajjeobimgagajmlibd}{Chrome
  extension for SSH}
\end{itemize}

\subsection{Connecting to CSD3}\label{connecting-to-csd3-1}

\begin{itemize}
\tightlist
\item
  host is \texttt{login.hpc.cam.ac.uk}
\item
  \href{https://docs.hpc.cam.ac.uk/hpc/user-guide/quickstart.html}{CSD3
  Docs}
\end{itemize}

\subsection{The Command Line}\label{the-command-line}

\begin{itemize}
\tightlist
\item
  Not as discoverable as a GUI
\item
  You can't break the HPC system
\item
  You type a command with optional flags and optional arguments and
  press ``Return''
\item
  The system may or may not give you any output
\end{itemize}

\subsection{The Command Line
resources}\label{the-command-line-resources}

\begin{itemize}
\tightlist
\item
  \href{https://swcarpentry.github.io/shell-novice/}{swcarpentry.github.io/shell-novice/}
\item
  \href{https://wizardzines.com/comics/every-core-unix-program-i-use/}{wizardzines.com/comics/every-core-unix-program-i-use/}
\end{itemize}

\subsection{The Scheduler}\label{the-scheduler}

The scheduler takes requests to run jobs with particular cluster
resources, fits these in around other user's jobs according to some
policy, launches the job, terminates the job if it is overrunning, does
accounting.

Examples:

\begin{itemize}
\tightlist
\item
  PBSpro
\item
  Platform LSF
\item
  Flux
\item
  Slurm (today, on CSD3)
\end{itemize}

\subsection{Job Scripts}\label{job-scripts}

A shell script with shell comments that are directives to the sheduler
about how the jobs should be run

\begin{Shaded}
\begin{Highlighting}[]
\CommentTok{\#!/bin/bash}
\CommentTok{\#SBATCH {-}{-}account=TRAINING{-}CPU }
\CommentTok{\#SBATCH {-}{-}reservation=iccs{-}summer{-}school1}
\CommentTok{\#SBATCH {-}{-}time=00:02:00}
\CommentTok{\#SBATCH {-}{-}job{-}name=my{-}first{-}job}
\CommentTok{\#SBATCH {-}{-}nodes=1}
\CommentTok{\#SBATCH {-}{-}cores=1}

\BuiltInTok{echo} \StringTok{"My first job {-} hooray"}
\end{Highlighting}
\end{Shaded}

\subsection{Submitting Jobs}\label{submitting-jobs}

\texttt{sbatch\ job.sh}

You will get back a Job ID.

\subsection{Viewing the Queue}\label{viewing-the-queue}

\begin{itemize}
\tightlist
\item
  \texttt{squeue}
\item
  \texttt{squeue\ -\/-me}
\end{itemize}

\subsection{Job Output}\label{job-output}

If you don't specify, by default it will be called
\texttt{slurm-\textless{}\$JOBID\textgreater{}.out}

To change this you can add an extra directive
\texttt{\#SBATCH\ -\/-output=}

\subsection{Exercise 2}\label{exercise-2}

\begin{itemize}
\tightlist
\item
  Write a job script to echo ``hello world''
\item
  Submit the job with \texttt{sbatch}
\item
  See it in the queue with \texttt{squeue\ -\/-me}
\item
  Find the output in te directory you submitted it from
  \texttt{ls\ -lrt}
\item
  Examine the output using \texttt{cat}
\end{itemize}

\subsection{Exercise 3}\label{exercise-3}

\begin{itemize}
\tightlist
\item
  add in the unix command \texttt{sleep\ 60}
\item
  find you job in the queue with \texttt{squeue\ -\/-me}
\item
  kill it with \texttt{scancel\ \textless{}JOBID\textgreater{}}
\end{itemize}

\subsection{Exercise 4}\label{exercise-4}

\begin{itemize}
\tightlist
\item
  change the sleep to 180 seconds
\item
  reduce the job request time to 1 minute
\item
  see what happens
\end{itemize}

\subsection{Environment Modules}\label{environment-modules}

\begin{itemize}
\tightlist
\item
  Multiple versions of the same software can be installed and you can
  choose between them
\item
  Two implementations

  \begin{itemize}
  \tightlist
  \item
    \href{https://modules.readthedocs.io/en/latest/}{Environment
    Modules} (TCL)
  \item
    \href{https://lmod.readthedocs.io/en/latest/}{LMod} (Lua)
  \end{itemize}
\item
  Module names are \textbf{site specific}
\item
  Module output is on \texttt{stderr} (!)

  \begin{itemize}
  \tightlist
  \item
    grep with
    \texttt{module\ avail\ 2\textgreater{}\&1\ \textbar{}\ grep\ the\_thing}
  \end{itemize}
\item
  If you load a module interactively remember to load it in a job script
\item
  Purge and only load the modules needed in the job script
\end{itemize}

\subsection{Exercise 5}\label{exercise-5}

\begin{itemize}
\tightlist
\item
  List the currently loaded modules \texttt{module\ list}
\item
  List the available modules \texttt{module\ avail}
\item
  Try running \texttt{h5perf\_serial}
\item
  Load the module \texttt{module\ load\ hdf5/1.12.1}
\item
  Now try \texttt{h5perf\_serial}
\end{itemize}

\subsection{Array Jobs}\label{array-jobs}

\begin{itemize}
\tightlist
\item
  map array member ID to problem domain, needn't necessarily be 1-dim
\item
  array members need to be independant can execute concurrently and
  order isn't guaranteed
\end{itemize}

\begin{Shaded}
\begin{Highlighting}[]
\CommentTok{\#!/bin/bash}
\CommentTok{\#SBATCH {-}{-}account=TRAINING{-}CPU }
\CommentTok{\#SBATCH {-}{-}reservation=iccs{-}summer{-}school2}
\CommentTok{\#SBATCH {-}{-}time=00:02:00}
\CommentTok{\#SBATCH {-}{-}job{-}name=array{-}test}
\CommentTok{\#SBATCH {-}{-}partition=icelake}
\CommentTok{\#SBATCH {-}{-}nodes=1}
\CommentTok{\#SBATCH {-}{-}cores=1}

\BuiltInTok{echo} \StringTok{"Hello from array member }\VariableTok{$\{SLURM\_ARRAY\_TASK\_ID\}}\StringTok{"}
\end{Highlighting}
\end{Shaded}

\texttt{sbatch\ -\/-array=1-10\ array-test.sh}

\subsection{Exercise 6}\label{exercise-6}

\begin{itemize}
\tightlist
\item
  submit an array job using the previous example
\item
  see how it appears in \texttt{squeue\ -\/-me}
\item
  examine the output files
\end{itemize}

\subsection{Workflows}\label{workflows}

\begin{itemize}
\item
  Do not be tempted to write your own workflow orchestrator
\item
  Choose from one of the already existing ones, e.g.:
\item
  \href{https://snakemake.readthedocs.io/en/stable/}{Snakemake}
\item
  \href{https://www.nextflow.io/docs/latest/}{NextFlow}
\end{itemize}

See \href{https://docs.nersc.gov/jobs/workflow/}{NERSC's advice}

\subsection{programming HPC}\label{programming-hpc}

\subsubsection{single node}\label{single-node}

\begin{itemize}
\tightlist
\item
  \href{https://www.openmp.org/}{OpenMP} is a specification for parallel
  programming
\item
  Hardware independent by design (e.g., CPU, FPGA, GPU\ldots)
\item
  \textbf{Shared memory} multiprocessing programming model
\end{itemize}

\begin{center}
\includegraphics[width=0.65\linewidth,height=\textheight,keepaspectratio]{imgs/openmp-logo.png}
\end{center}

\begin{center}\rule{0.5\linewidth}{0.5pt}\end{center}

\subsubsection{hello OpenMP}\label{hello-openmp}

\begin{Shaded}
\begin{Highlighting}[]
\PreprocessorTok{\#include }\ImportTok{\textless{}stdio.h\textgreater{}}
\PreprocessorTok{\#include }\ImportTok{\textless{}omp.h\textgreater{}}

\DataTypeTok{int}\NormalTok{ main}\OperatorTok{()} \OperatorTok{\{}
  \PreprocessorTok{\#pragma omp parallel}
  \OperatorTok{\{}
    \DataTypeTok{int}\NormalTok{ thread\_id }\OperatorTok{=}\NormalTok{ omp\_get\_thread\_num}\OperatorTok{();}
\NormalTok{    printf}\OperatorTok{(}\StringTok{"Hello from thread }\SpecialCharTok{\%d\textbackslash{}n}\StringTok{"}\OperatorTok{,}\NormalTok{ thread\_id}\OperatorTok{);}
  \OperatorTok{\}}
  \ControlFlowTok{return} \DecValTok{0}\OperatorTok{;}
\OperatorTok{\}}
\end{Highlighting}
\end{Shaded}

Build with

\begin{Shaded}
\begin{Highlighting}[]
\ExtensionTok{$}\NormalTok{ gcc }\AttributeTok{{-}fopenmp} \AttributeTok{{-}O3}\NormalTok{ hello.c }\AttributeTok{{-}o}\NormalTok{ hello.exe}
\end{Highlighting}
\end{Shaded}

\begin{center}\rule{0.5\linewidth}{0.5pt}\end{center}

\subsubsection{hello OpenMP}\label{hello-openmp-1}

\begin{itemize}
\tightlist
\item
  Can you predict what the output will be with 4 threads?
\end{itemize}

\begin{tcolorbox}[enhanced jigsaw, arc=.35mm, bottomrule=.15mm, bottomtitle=1mm, breakable, colback=white, colbacktitle=quarto-callout-note-color!10!white, colframe=quarto-callout-note-color-frame, coltitle=black, left=2mm, leftrule=.75mm, opacityback=0, opacitybacktitle=0.6, rightrule=.15mm, title=\textcolor{quarto-callout-note-color}{\faInfo}\hspace{0.5em}{Note}, titlerule=0mm, toprule=.15mm, toptitle=1mm]

To change the number of threads we need to set environment variable
\texttt{OMP\_NUM\_THREADS} e.g., \texttt{export\ OMP\_NUM\_THREADS=1}

\end{tcolorbox}

\begin{center}\rule{0.5\linewidth}{0.5pt}\end{center}

\subsubsection{Exercise 7}\label{exercise-7}

\begin{itemize}
\tightlist
\item
  Build the OpenMP hello world
\item
  Run with varying numbers of threads
\end{itemize}

\subsubsection{hello OpenMP}\label{hello-openmp-2}

What happens when we run the following code?

\begin{Shaded}
\begin{Highlighting}[]
\ExtensionTok{$}\NormalTok{ gcc }\AttributeTok{{-}fopenmp} \AttributeTok{{-}O3}\NormalTok{ hello.c }\AttributeTok{{-}o}\NormalTok{ hello.exe}
\ExtensionTok{$}\NormalTok{ ./hello.exe}
\ExtensionTok{hello}\NormalTok{ from process 1}
\ExtensionTok{hello}\NormalTok{ from process 3}
\ExtensionTok{hello}\NormalTok{ from process 0}
\ExtensionTok{hello}\NormalTok{ from process 2}
\end{Highlighting}
\end{Shaded}

\begin{center}\rule{0.5\linewidth}{0.5pt}\end{center}

\subsubsection{OpenMP Scaling}\label{openmp-scaling}

\begin{itemize}
\tightlist
\item
  Try running the \texttt{pi.c} example (in directory
  \texttt{example-code/})
\item
  To compile \texttt{pi.c} use command:
\end{itemize}

\begin{Shaded}
\begin{Highlighting}[]
\BuiltInTok{cd}\NormalTok{ example{-}code}
\FunctionTok{gcc} \AttributeTok{{-}fopenmp} \AttributeTok{{-}O3}\NormalTok{ pi.c }\AttributeTok{{-}o}\NormalTok{ pi.exe}
\end{Highlighting}
\end{Shaded}

\begin{center}\rule{0.5\linewidth}{0.5pt}\end{center}

\subsubsection{What might be causing
this?}\label{what-might-be-causing-this}

\begin{itemize}
\tightlist
\item
  The problem is ``pleasingly parallel''
\item
  Process pinning
\item
  NUMA

  \begin{itemize}
  \tightlist
  \item
    Non-unifrom memory architecture
  \item
    Differing latencies from cores to main memory
  \end{itemize}
\item
  First touch placement policy.
\end{itemize}

\begin{center}\rule{0.5\linewidth}{0.5pt}\end{center}

\subsubsection{ls-topo}\label{ls-topo}

\pandocbounded{\includegraphics[keepaspectratio]{imgs/lstopo-login-q.png}}

\subsection{programming HPC}\label{programming-hpc-1}

\subsubsection{distributed (multi-node)}\label{distributed-multi-node}

\begin{itemize}
\tightlist
\item
  MPI (Message Passing Interface)
\item
  One of the most common methods of distributed compute
\item
  \textbf{Distributed memory} multiprocessing programming model
\item
  Implementations (openMPI, MPICH, Intel MPI)
\end{itemize}

\begin{center}\rule{0.5\linewidth}{0.5pt}\end{center}

\subsubsection{MPI Hello World}\label{mpi-hello-world}

\begin{Shaded}
\begin{Highlighting}[]
\PreprocessorTok{\#include }\ImportTok{\textless{}mpi.h\textgreater{}}
\PreprocessorTok{\#include }\ImportTok{\textless{}stdio.h\textgreater{}}

\DataTypeTok{int}\NormalTok{ main}\OperatorTok{(}\DataTypeTok{int}\NormalTok{ argc}\OperatorTok{,} \DataTypeTok{char}\OperatorTok{**}\NormalTok{ argv}\OperatorTok{)} \OperatorTok{\{}
    \CommentTok{// Initialize the MPI environment}
\NormalTok{    MPI\_Init}\OperatorTok{(}\NormalTok{argc}\OperatorTok{,}\NormalTok{ argv}\OperatorTok{);}

    \CommentTok{// Get the number of processes}
    \DataTypeTok{int}\NormalTok{ world\_size}\OperatorTok{;}
\NormalTok{    MPI\_Comm\_size}\OperatorTok{(}\NormalTok{MPI\_COMM\_WORLD}\OperatorTok{,} \OperatorTok{\&}\NormalTok{world\_size}\OperatorTok{);}

    \CommentTok{// Get the rank of the process}
    \DataTypeTok{int}\NormalTok{ world\_rank}\OperatorTok{;}
\NormalTok{    MPI\_Comm\_rank}\OperatorTok{(}\NormalTok{MPI\_COMM\_WORLD}\OperatorTok{,} \OperatorTok{\&}\NormalTok{world\_rank}\OperatorTok{);}

    \CommentTok{// Get the name of the processor}
    \DataTypeTok{char}\NormalTok{ processor\_name}\OperatorTok{[}\NormalTok{MPI\_MAX\_PROCESSOR\_NAME}\OperatorTok{];}
    \DataTypeTok{int}\NormalTok{ name\_len}\OperatorTok{;}
\NormalTok{    MPI\_Get\_processor\_name}\OperatorTok{(}\NormalTok{processor\_name}\OperatorTok{,} \OperatorTok{\&}\NormalTok{name\_len}\OperatorTok{);}

    \CommentTok{// Print off a hello world message}
\NormalTok{    printf}\OperatorTok{(}\StringTok{"Hello world from processor }\SpecialCharTok{\%s}\StringTok{, rank }\SpecialCharTok{\%d}\StringTok{ out of }\SpecialCharTok{\%d}\StringTok{ processors}\SpecialCharTok{\textbackslash{}n}\StringTok{"}\OperatorTok{,}
\NormalTok{           processor\_name}\OperatorTok{,}\NormalTok{ world\_rank}\OperatorTok{,}\NormalTok{ world\_size}\OperatorTok{);}

    \CommentTok{// Finalize the MPI environment.}
\NormalTok{    MPI\_Finalize}\OperatorTok{();}
\OperatorTok{\}} 
\end{Highlighting}
\end{Shaded}

\begin{center}\rule{0.5\linewidth}{0.5pt}\end{center}

\subsubsection{Building MPI Programs}\label{building-mpi-programs}

Build:

\begin{Shaded}
\begin{Highlighting}[]
\ExtensionTok{mpicc}\NormalTok{ mpi{-}hello.c }\AttributeTok{{-}o}\NormalTok{ mpi{-}hello.exe}
\end{Highlighting}
\end{Shaded}

Run a small test on login node:

\begin{verbatim}
mpiexec -np 4 ./mpi-hello.exe
\end{verbatim}

\subsubsection{Exercise 8}\label{exercise-8}

\begin{itemize}
\tightlist
\item
  Build \texttt{mpi-hello.c}
\item
  Write a job script to run on 2 nodes, 3 cores per node

  \begin{itemize}
  \tightlist
  \item
    (\texttt{-np} would be 2*3 = 6)
  \end{itemize}
\end{itemize}

\subsection{programming HPC}\label{programming-hpc-2}

\subsubsection{GPU offloading}\label{gpu-offloading}

\begin{itemize}
\tightlist
\item
  There are a range of GPU offloading programming models
\item
  Vendor specific

  \begin{itemize}
  \tightlist
  \item
    \href{https://developer.nvidia.com/cuda-toolkit}{CUDA}
  \item
    \href{https://rocm.docs.amd.com/projects/HIP/en/docs-develop/what_is_hip.html}{HIP}
  \end{itemize}
\item
  Vendor agnostic

  \begin{itemize}
  \tightlist
  \item
    \href{https://www.khronos.org/sycl/}{SYCL}
  \item
    \href{https://kokkos.org/kokkos-core-wiki/}{kokkos}
  \item
    openMP GPU offloading
  \item
    \href{https://www.openacc.org/}{openAcc}
  \end{itemize}
\end{itemize}

\subsection{The bad news: Amdahal's
law}\label{the-bad-news-amdahals-law}

\(S = 1 / (1 - p + p/s)\), where\ldots{}

\begin{itemize}
\tightlist
\item
  \(S\) is the speedup of a process
\item
  \(p\) is the proportion of a program that can be made parallel, and
\item
  \(s\) is the speedup factor of the parallel portion.
\end{itemize}

\subsection{Amdahal's Law}\label{amdahals-law}

\pandocbounded{\includegraphics[keepaspectratio]{imgs/AmdahlsLaw.svg.png}}

\subsection{The good news: Gustafson's
law}\label{the-good-news-gustafsons-law}

\begin{itemize}
\tightlist
\item
  Mitigates the drawbacks of Amdahl's Law
\item
  Scale the problem size as you scale the number of nodes
\end{itemize}

\subsection{debugging}\label{debugging}

\begin{itemize}
\tightlist
\item
  Multiple strategies {🔍}{🐛}

  \begin{itemize}
  \tightlist
  \item
    \texttt{printf()}
  \item
    logging
  \item
    debuggers (\texttt{gdb}, \texttt{lldb}, linaro \texttt{ddt}\ldots)
  \end{itemize}
\item
  \href{https://www.sourceware.org/gdb/}{gdb}

  \begin{itemize}
  \tightlist
  \item
    available on most HPC systems
  \item
    works with C, C++, Fortran, Rust\ldots{}
  \item
    Command-line interface
  \end{itemize}
\item
  \href{https://github.com/Cambridge-ICCS/debugging-workshop}{Debugging
  Course} from Summer School 2025
  (\href{https://youtu.be/Oo3cPIckI0Y?si=tFr1UYZLtP8vkppH}{video}
\end{itemize}

\section{Profiling}\label{profiling}

\begin{tcolorbox}[enhanced jigsaw, arc=.35mm, bottomrule=.15mm, bottomtitle=1mm, breakable, colback=white, colbacktitle=quarto-callout-warning-color!10!white, colframe=quarto-callout-warning-color-frame, coltitle=black, left=2mm, leftrule=.75mm, opacityback=0, opacitybacktitle=0.6, rightrule=.15mm, title=\textcolor{quarto-callout-warning-color}{\faExclamationTriangle}\hspace{0.5em}{Warning!}, titlerule=0mm, toprule=.15mm, toptitle=1mm]

Premature Optimization Is the Root of All Evil

\end{tcolorbox}

Donald Knuth (1974)

\subsection{Profiling}\label{profiling-1}

\begin{figure}[H]

{\centering \includegraphics[width=0.5\linewidth,height=\textheight,keepaspectratio]{imgs/rust-dev-survey-profiling.png}

}

\caption{\href{https://www.jetbrains.com/lp/devecosystem-2023/rust/}{JetBrains
rust developer survey}}

\end{figure}%

\subsection{Profiling}\label{profiling-2}

\begin{itemize}
\tightlist
\item
  \href{https://www.intel.com/content/www/us/en/docs/vtune-profiler/get-started-guide/2025-0/windows-os.html}{Intel
  VTune}
\item
  \href{https://docs.olcf.ornl.gov/software/profiling/Scorep.html}{Score-P}
\item
  \href{https://www.cs.uoregon.edu/research/tau/about.php}{TAU}
\item
  \href{https://hpctoolkit.org/}{HPCToolkit}
\item
  \href{https://github.com/RRZE-HPC/likwid?tab=readme-ov-file}{LIKWID}
\end{itemize}

\subsection{Example with Score-p}\label{example-with-score-p}

\begin{itemize}
\tightlist
\item
  We introduce a small simple OpenMP
  \href{https://github.com/niravshah241/cpp_profiler_example/tree/matrix_vector_multiplication}{example}.
  Precisely, we introduce dense-Matrix vector multiplication example,
  where we implement element wise multiplication.
\item
  We consider parallelization across row indices using multiple threads.
\item
  Compilation step (Notice the use of \texttt{scorep-g++} instead of
  \texttt{g++})
\end{itemize}

\begin{Shaded}
\begin{Highlighting}[]
\ExtensionTok{scorep{-}g++} \AttributeTok{{-}c} \AttributeTok{{-}fopenmp}\NormalTok{ main.cpp generate\_random\_matrix.cpp matrix\_multiply.cpp}
\end{Highlighting}
\end{Shaded}

\begin{itemize}
\tightlist
\item
  Linking step (Notice the use of \texttt{scorep-g++} instead of
  \textt{g++})
\end{itemize}

\begin{Shaded}
\begin{Highlighting}[]
\ExtensionTok{scorep{-}g++} \AttributeTok{{-}o}\NormalTok{ main\_executable }\AttributeTok{{-}fopenmp}\NormalTok{ main.o generate\_random\_matrix.o matrix\_multiply.o}
\end{Highlighting}
\end{Shaded}

\begin{itemize}
\tightlist
\item
  Program Execution
\end{itemize}

\begin{Shaded}
\begin{Highlighting}[]
\ExtensionTok{./main\_executable}
\end{Highlighting}
\end{Shaded}

\begin{itemize}
\tightlist
\item
  A folder containing profiling data using \texttt{.cubex} and
  \texttt{.cfg} file is generated.
\end{itemize}

\subsection{IO profiling with Darshan}\label{io-profiling-with-darshan}

\subsection{Green HPC}\label{green-hpc}

\begin{itemize}
\tightlist
\item
  Green500 - Top500 divided by energy
\item
  Calculators:

  \begin{itemize}
  \tightlist
  \item
    \href{https://www.green-algorithms.org/}{Green Algorithms}
  \item
    \href{https://chryswoods.github.io/howmuchisenough/}{How Much is
    Enough}
  \end{itemize}
\item
  The Carpentries
  \href{https://carpentries-incubator.github.io/green-software-hpc/index.html}{Green
  software use on HPC}
\end{itemize}

Advice:

\begin{itemize}
\tightlist
\item
  Profile!
\item
  Use reduced precision
\item
  Design experiments well
\end{itemize}

\subsection{Exercise}\label{exercise}

\begin{itemize}
\tightlist
\item
  Go to the \href{https://www.green-algorithms.org/}{Green Algorithms}
  website
\item
  Use the calculator to look at

  \begin{itemize}
  \tightlist
  \item
    3 hours
  \item
    4 nodes (4 * 76 cores)
  \item
    Xeon Platinum 8268
  \item
    Memory 512 GB
  \item
    Location Europe / United Kingdom
  \item
    PUE 1.1
  \end{itemize}
\end{itemize}

\subsection{Applying for Resources}\label{applying-for-resources}

\begin{itemize}
\tightlist
\item
  \href{https://doeleadershipcomputing.org/}{Incite}
\item
  \href{https://eurohpc-ju.europa.eu/supercomputers/supercomputers-access-calls_en}{Euro
  Joint Undertaking}
\end{itemize}

\subsection{Further Resources}\label{further-resources}

\begin{itemize}
\tightlist
\item
  Your local HPC support
\item
  HPC carpentry

  \begin{itemize}
  \tightlist
  \item
    https://carpentries-incubator.github.io/hpc-intro/
  \end{itemize}
\item
  ARCHER2

  \begin{itemize}
  \tightlist
  \item
    https://www.archer2.ac.uk/training/materials/
  \end{itemize}
\item
  ATPESC

  \begin{itemize}
  \tightlist
  \item
    https://extremecomputingtraining.anl.gov/
  \end{itemize}
\item
  SC, ISC Tutorials
\end{itemize}

\subsection{Contact}\label{contact}

For more information we can be reached at:

\faIcon{pencil} ~Chris Edsall

\faIcon{person-digging}
~\href{https://iccs.cam.ac.uk/about-us/our-team}{ICCS/UoCambridge}

\faIcon{envelope} ~\href{mailto:cje57@cam.ac.uk}{cje57{[}AT{]}cam.ac.uk}

\faIcon{github}
~\href{https://github.com/christopheredsall}{christopheredsall}

\faIcon{mastodon}
~\href{https://scholar.social/@hpcchris}{@hpcchris@scholar.social}

You can also contact the ICCS,
\href{https://iccs.cam.ac.uk/resources-vesri-members/resource-allocation-process}{make
a resource allocation request}, or visit us at the
\href{https://docs.google.com/spreadsheets/d/1v8y2GodI9JZoHrFRpLW2tD135MurDBMIcr-ArBFIT3A}{Summer
School RSE Helpdesk}.




\end{document}
